\documentclass[11pt]{article}

\usepackage[a4paper,margin=2.5cm]{geometry}

\usepackage{amsmath,amssymb}
\usepackage{microtype}

\usepackage[T1]{fontenc}
% \usepackage{lmodern}   % lmodern.sty absent on this TeX install; falls back to Computer Modern. Re-enable once installed.

\usepackage{graphicx}
\usepackage{xcolor}
\usepackage{booktabs}
\usepackage{enumitem}
\usepackage{caption}
\usepackage{hyperref}

% ------------------------------------------------------------------
% Nature / Cell review style
% ------------------------------------------------------------------

% No paragraph indentation anywhere
\setlength{\parindent}{0pt}

% Space between paragraphs
\setlength{\parskip}{0.75em}

% Slightly looser line spacing
\linespread{1.05}

% Better page breaking
\raggedbottom

% Slightly larger display math spacing
\setlength{\abovedisplayskip}{1.0em}
\setlength{\belowdisplayskip}{1.0em}
\setlength{\abovedisplayshortskip}{0.8em}
\setlength{\belowdisplayshortskip}{0.8em}

% Compact lists
\setlist{
    itemsep=0.3em,
    topsep=0.4em
}

% Figure captions
\captionsetup{
    font=small,
    labelfont=bf
}

% Hyperlinks
\hypersetup{
    colorlinks=true,
    linkcolor=blue,
    citecolor=blue,
    urlcolor=blue
}

% Convenience macros
\newcommand{\code}[1]{\texttt{#1}}

% After displayed equations, continue flush-left
\newcommand{\aftereq}{\par\noindent}

\title{Knowledge-First Inverse Modelling:\\ Mechanism Discovery over Parameter Optimization\\[4pt]
\large A cardiomyocyte active-stress MPM case study in Plexus}
\author{}
\date{}

\begin{document}
\maketitle

\begin{abstract}
Modern differentiable simulators make parameter optimization easy. They do not make mechanism discovery easy. In many inverse-modelling problems the natural objective is to find the parameter vector that best fits the data. This parameter-optimization view treats the simulator as a black box: propose parameters, run the model, compute a score, move toward a better score. It is efficient when the forward model is already correct, but it becomes misleading when the model is incomplete: a better score may only mean the optimizer has found a compensating artifact, not that the model has captured the mechanism. A knowledge-first strategy uses optimization differently. Each batch is not only a step toward a better fit; it is an experiment designed to learn how the system works. Parameters are mechanistic levers, a failed run is useful when it falsifies a hypothesis, and a successful run is useful only when its improvement can be explained in terms of morphology, dynamics, or physical coupling. The objective is therefore not just $\theta^* = \arg\min_\theta \mathcal{L}(\theta)$ but the construction of a map
\[
\theta \;\longrightarrow\; \text{mechanism} \;\longrightarrow\; \text{morphology} \;\longrightarrow\; \text{fit}.
\]
We make this concrete on the inverse problem of fitting real cardiomyocyte tissue motion with a differentiable active-stress material-point-method (MPM) model built from registered Plexus operators, scored by loop \emph{morphology} (the LoopScore metric) rather than framewise displacement. After many batches the deliverable is not a single optimum but a set of mechanistic constraints on what the forward model must look like: \textbf{each mechanistic insight removes an entire class of incorrect models}, and a fit that plateaus despite continued optimization is read not as an optimization failure but as evidence that the forward model is missing a genuine physical ingredient. This closes a discovery loop that we take to be a defining move of the method~--- \emph{optimization discovers residuals; residuals generate hypotheses; hypotheses become operators; operators extend the language}:
\[
\text{residual} \;\longrightarrow\; \text{hypothesis} \;\longrightarrow\; \text{operator} \;\longrightarrow\; \text{extended language}.
\]
A residual the \emph{current} operator language cannot break is therefore not a dead end but a prompt: the accumulated knowledge names the missing physical ingredient, and it is added as a new registered operator. The central claim thus has a constructive counterpart~--- the inverse problem becomes an engine for discovering the operators a correct forward model requires, so that fitting the data and \emph{building the language that can express it} are the same activity.
\end{abstract}

\section{Introduction}

Each batch is not only a step toward a better fit; it is an experiment designed to learn how the system works. Parameters are mechanistic levers, a failed run is useful when it falsifies a hypothesis, and a successful run is useful only when its improvement can be explained in terms of morphology, dynamics, or physical coupling. The objective is therefore not just $\theta^* = \arg\min_\theta \mathcal{L}(\theta)$ but the construction of a map
\[
\theta \;\longrightarrow\; \text{mechanism} \;\longrightarrow\; \text{morphology} \;\longrightarrow\; \text{fit}.
\]
We make this concrete on the inverse problem of fitting real cardiomyocyte tissue motion with a differentiable active-stress material-point-method (MPM) model built from registered Plexus operators, scored by loop \emph{morphology} (the LoopScore metric) rather than framewise displacement. After many batches the deliverable is not a single optimum but a set of mechanistic constraints on what the forward model must look like. The central claim: \textbf{each mechanistic insight removes an entire class of incorrect models}, and a fit that plateaus despite continued optimization is read not as an optimization failure but as evidence that the forward model is missing a genuine physical ingredient~--- with the accumulated knowledge naming which ingredient to add next.

This distinction matters most exactly when optimization stalls. When the score stops improving despite continued training and small architectural changes, the parameter-optimization view sees a disappointing plateau; the knowledge-first view reads the \emph{same} plateau as a measurement~--- evidence that the forward model is missing a genuine physical ingredient, not that the optimizer has failed. That re-reading is only available if each run has been treated as an experiment: the parameters must be mechanistic levers whose effect on morphology is recorded, the score must measure the object of interest (here loop \emph{shape}, not framewise displacement), and the accumulated knowledge must be organised so that a residual can be attributed to a specific morphological axis rather than a generic ``too small.''

When a residual survives every composition of the operators already in the model, that irreducibility is itself the result. It licenses a specific move: \emph{extend the language.} Optimization discovers a residual; the residual, attributed to a morphological axis, generates a hypothesis about the missing physics; the hypothesis is realised as a new registered operator; and the operator enlarges the vocabulary the forward model can express,
\[
\text{residual} \;\longrightarrow\; \text{hypothesis} \;\longrightarrow\; \text{operator} \;\longrightarrow\; \text{extended language}.
\]
Fitting the data and \emph{building the language that can express it} thereby become the same activity. In the case study below this loop runs to completion more than once. A plateau in loop \emph{enclosure}~--- trajectories that trace thin, nearly one-dimensional arcs rather than closed loops~--- is traced not to insufficient force but to a symmetry: a contraction axis that is fixed over the beat retraces its path and encloses no area. The named ingredient is a rotating contraction axis, added as a single operator, which breaks the plateau and fills the enclosure residual. A subsequent plateau in loop \emph{size} then resists every lever the current language offers~--- drive, gain, stiffness, duration~--- and prompts a different hypothesis: that the tissue does not enter each beat from a stress-free reference, so that the missing ingredient is a residual-stress (pre-stress) operator, proposed precisely because no operator already in the model can reach it. Each such step both removes a class of models and adds a word to the vocabulary.

\section{From a continuum with a pattern to a tissue made of cells}

The campaign so far fits a \emph{continuum}. Heterogeneity enters as a map: an image is sampled at
each material point and becomes that point's Young's modulus, so neighbouring points differ a
little and the contraction's symmetry is broken. This is enough to produce spatial variation and it
is not enough to produce a tissue, because a map has no boundaries in it. Two material points a
micron apart always have almost the same material, whatever the map. A monolayer of cardiomyocytes
does not work that way: it is a set of discrete contractile units, mechanically coupled but
individually stiff, individually oriented, and individually timed. The residual this leaves is
structural rather than numerical --- no setting of a smooth field can make two adjacent points
belong to different objects --- and by the argument of this paper it is exactly the kind of
residual that names its own missing ingredient.

\paragraph{The hypothesis comes from the data, and from its motion rather than its appearance.}
Cell boundaries are not visible in these recordings. The phase-contrast frames are a dense texture
at confluence, and segmenting them is segmenting nothing --- an edge detector finds halos and
granules. But a cardiomyocyte contracts along its own long axis, and that axis is a property of the
cell, not of the image. It can therefore be measured \emph{over time at one point} before any
spatial reasoning: each grid point's beat trajectory is nearly a straight line (median anisotropy
$\lambda_1/\lambda_2 = 14$) whose direction is that point's contraction axis. Segmentation then
becomes: cut where the axis turns, using $\exp(2i\theta)$ because an axis has no head or tail.

\paragraph{What that measurement does and does not establish.} It is worth stating both, because
the positive results are all of one kind. The axis field is real rather than an artefact of
smoothing: computed from two disjoint pairs of beats it agrees at $0.9989$ with \emph{no} smoothing,
and its spatial coherence length is $90$\,px. The resulting partition reproduces across beats at
$\mathrm{ARI} = 0.95$ against a rolled null of $0.18$, and it recovers the right cell density
without being told it --- $465$ regions against $472$ nuclei counted independently in the image, a
$1.5\%$ agreement, which also pins the uncalibrated pixel size at $\approx 0.33\,\mu$m so that the
$90$\,px coherence length is $30\,\mu$m, a cardiomyocyte.

Against that: the unseeded partition does \emph{not} localise individual cells. Displace the nuclei
and they fall one-per-region just as often ($z = +0.5$), and the borders do not coincide with image
edges ($z = -1.9$). Reproducible is not correct: a patch of aligned myofibrils produces the same
stable structure and may span two cells or divide one. The construction that survives its nulls
uses the nuclei --- which are visible --- to fix \emph{where} the cells are, and the motion to
decide only \emph{where the border between two of them runs}. Scored by fitting one affine map per
region to held-out beats, on the principle that a cell is the unit that moves together, the
motion-bent partition explains the held-out motion better than a Voronoi tessellation of the same
nuclei by $2.1\%$, against $1.1\%$ for the same borders displaced --- $z = +3.1$, and $+3.2$ with
the halves exchanged. Small, and earned.

\paragraph{The operator.} The ingredient is then named and registered:
\code{seed\_from\_segmentation}, with a companion \code{label\_image} field. Two properties of a
label map force both to exist. Labels may not be interpolated --- cell $7$ beside cell $12$ must not
yield cell $9$ in between --- so the field is read without normalisation and sampled by nearest
neighbour, unlike the greyscale map field it otherwise resembles. And a cell is not a disc: the
engine scatters a contained set in a ball about its parent, which is correct for a generic child
and wrong for a cell whose shape has been measured, so the operator re-seeds each cell entity onto
its own centroid and places that cell's material points inside its own mask. It is structural, not
dynamics: it establishes a configuration and then returns immediately, and it declares the
exemption from the integration invariant that writing state directly requires.

\paragraph{What the language can now say.} The hierarchy the framework describes --- sets contained
in sets --- becomes populated at three levels that each mean something:
\[
\text{tissue} \;\longrightarrow\; \text{cells} \;\longrightarrow\; \text{material points},
\]
one tissue, $472$ cells measured from the recording, and $500$ material points per cell placed
inside its measured shape. Every point of a given cell carries that cell's Young's modulus
\emph{exactly}, so two points either side of a boundary differ discontinuously --- which is what a
painted map cannot do, and which lets adjacent cells shear against one another. The per-cell
stiffness is not invented either: it is taken from each cell's measured beat amplitude, under the
stated hypothesis that a cell which moved little is stiff rather than weakly contractile. That
hypothesis is a lever, and the alternative reading --- amplitude as contraction gain --- is the same
line the other way round, which makes it a testable fork rather than a modelling choice buried in
code.

This is one full turn of the loop the paper argues for. A residual that no parameter could reach
(the model has no cells) generated a hypothesis (the tissue is a set of discrete units), the
hypothesis was measured in the data by a route the appearance of the data forbids (time before
space), and the result became a registered operator that enlarges what the forward model can
express. Whether it improves the \emph{fit} is the next experiment, and it is now a question the
language can pose.

The remainder of the paper develops this method and its instantiation: the differentiable active-stress MPM forward model and the operators it is assembled from; the LoopScore morphology metric and why framewise error is the wrong objective for a loop; the residual decomposition that attributes a plateau to a named morphological axis; and the sequence of operator-discovery steps the campaign has produced, presented as the running example of optimization, residual, hypothesis, operator, extended language.

\end{document}